% !TeX root = ../main.tex
\chapter{Agregados: aceitação de doação e ciclo de vida}
\label{ap:agregados}

Este apêndice reúne trechos das raízes de agregado \codigo{DonationRequest} e \codigo{Donation}, evidenciando a proteção de invariantes no núcleo do domínio --- princípio tático do DDD discutido no Capítulo~2 e materializado na modelagem do Capítulo~5.

\section{Aceitação de doação pela solicitação}

A Listagem~\ref{lst:accepts-donation} mostra o método \codigo{acceptsDonation}, usado pelo algoritmo de cumprimento de metas (Apêndice~\ref{ap:fulfillment}). A solicitação só aceita a doação se estiver ativa, não expirada, se a doação estiver concluída, se a data da doação estiver na janela \texttt{[dateRequested, dateLimit]} e se o tipo sanguíneo do doador for compatível com o tipo solicitado.

\lstinputlisting[
  language=Java,
  caption={Método \texttt{DonationRequest.acceptsDonation}},
  label={lst:accepts-donation}
]{apendices/codigo/DonationRequestAcceptsDonation.java}
\fonte{Elaborado pelos próprios autores (2026).}

\section{Conclusão de doação pendente}

A Listagem~\ref{lst:donation-complete} ilustra a máquina de estados da doação: \codigo{validateStatusFlags} exige exatamente um estado ativo entre pendente, concluída e cancelada; \codigo{complete} só permite a transição a partir de pendente e rejeita datas futuras.

\lstinputlisting[
  language=Java,
  caption={Validação de estado e método \texttt{Donation.complete}},
  label={lst:donation-complete}
]{apendices/codigo/DonationComplete.java}
\fonte{Elaborado pelos próprios autores (2026).}
